% !PRINTED 471
\weberpage{471}
%<<CONT

\gesperrt{2. Von den durch die symbolische Gleichung $\sigma \sigma' = \sigma''$ verbundenen drei Substitutionen des Körpers $\Omega(\varrho)$ können irgend zwei beliebig gegeben sein, während die dritte dadurch eindeutig bestimmt ist.}

Wir haben die Composition der Substitutionen symbolisch durch das Zeichen der Multiplication ausgedrückt. Es ist hier aber wohl auf die Reihenfolge der Factoren oder Componenten zu achten, weil $\sigma \sigma'$ von $\sigma' \sigma$ verschieden sein kann; d.\,h. bei der die Composition ausdrückenden symbolischen Multiplication gilt nicht das commutative Gesetz der gewöhnlichen Multiplication; wohl aber gilt das sogenannte associative Gesetz, das sich in folgendem Satz ausspricht:

\gesperrt{3. Sind $\sigma, \sigma', \sigma''$ irgend drei der $\mu$ Substitutionen des Körpers $\Omega(\varrho)$, so ist}
\weq{3}{(\sigma \sigma') \sigma'' = \sigma (\sigma' \sigma'').}
Der Beweis ist sehr einfach; denn nach dem Satz 3., \S.\ 146 können wir $\varrho_a, \varrho_b, \varrho_c$ so bestimmen, dass
\weq{4}{\sigma = (\varrho, \varrho_a), \quad \sigma' = (\varrho_a, \varrho_b), \quad \sigma'' = (\varrho_b, \varrho_c),}
und dann ist
\[
\sigma \sigma' = (\varrho, \varrho_b), \qquad (\sigma \sigma') \sigma'' = (\varrho, \varrho_b)(\varrho_b, \varrho_c) = (\varrho, \varrho_c),
\]
\[
\sigma' \sigma'' = (\varrho_a, \varrho_c), \qquad \sigma (\sigma' \sigma'') = (\varrho, \varrho_a)(\varrho_a, \varrho_c) = (\varrho, \varrho_c).
\]
Wir bezeichnen daher kurz die aus den drei Substitutionen $\sigma, \sigma', \sigma''$ zusammengesetzte Substitution mit
\[
\sigma \sigma' \sigma'',
\]
und können überhaupt aus beliebig vielen Componenten eine bestimmte Substitution
\[
\sigma \sigma' \sigma'' \sigma''' \ldots
\]
dadurch zusammensetzen, dass wir nach einander immer zwei benachbarte unter den Componenten zu einer einzigen Substitution vereinigen, bis das ganze symbolische Product sich auf eine einzige Substitution zusammengezogen hat. Denn wir können nach \S.\ 146, 3.\ setzen
\[
\sigma = (\varrho, \varrho_a), \quad \sigma' = (\varrho_a, \varrho_b), \quad \sigma'' = (\varrho_b, \varrho_c), \quad \sigma''' = (\varrho_c, \varrho_e),
\]
und erhalten immer
\[
(\varrho, \varrho_a)(\varrho_a, \varrho_b)(\varrho_b, \varrho_c)(\varrho_c, \varrho_e) = (\varrho, \varrho_e),
\]
was offenbar auf eine beliebige Anzahl von Componenten ausgedehnt werden kann.
